% from ega3-6-fr.tex
\oldpage[III]{150}
dire qu'ils sont de dimension cohomologique $\leqslant n$ signifie qu'ils sont
réguliers et de dimension (de Krull) $\leqslant n$ ($0_{\mathrm{IV}}$, 17.3.1).
Avec la terminologie de la théorie de la dimension que nous introduirons au
chap. IV, il revient au même de dire qu'un préschéma localement noethérien $T$
est de dimension cohomologique $\leqslant n$, ou de dire qu'il est régulier
($0_{\mathrm I}$, 4.1.4) et de dimension $\leqslant n$; cela signifie que pour
tout ouvert affine $U$ de $T$, l'anneau $\Gamma(U,\mathscr O_T)$ est de
dimension cohomologique $\leqslant n$ ($0_{\mathrm{IV}}$, 17.2.6). Cela étant,
cette dernière remarque, jointe à (6.3.2), prouve que si $S$ est localement
noethérien et de dimension cohomologique finie, les suites spectrales
$'\mathscr E(\mathscr P_\bullet,\mathscr Q_\bullet)$ et
$''\mathscr E(\mathscr P_\bullet,\mathscr Q_\bullet)$ sont birégulières.

Il est clair que
$\mathscr{T}\!or_\bullet^S(\mathscr P_\bullet,\mathscr Q_\bullet)$ se
transforme en
$\mathscr{T}\!or_\bullet^S(\mathscr Q_\bullet,\mathscr P_\bullet)$ (à un
isomorphisme près) par l'isomorphisme canonique de $X\times_S Y$ sur
$Y\times_S X$.
\end{env}

\begin{proposition}[6.5.6]
Soit $(\mathscr P_{\alpha\bullet})$ un système inductif filtrant de complexes
de $\mathscr O_X$-Modules quasi-cohérents; il existe alors un isomorphisme
canonique
\[
  \varinjlim_\alpha
  \mathscr{T}\!or_\bullet^S
  (\mathscr P_{\alpha\bullet},\mathscr Q_\bullet)
  \xrightarrow{\sim}
  \mathscr{T}\!or_\bullet^S
  \left(\varinjlim_\alpha\mathscr P_{\alpha\bullet},\mathscr Q_\bullet\right).
  \tag{6.5.6.1}\label{III.6.5.6.1.fr}
\]
\end{proposition}

\begin{proof}
La question étant locale sur $S$, $X$ et $Y$, on peut supposer $S$, $X$, $Y$
affines et la proposition se réduit alors à (6.3.6).
\end{proof}

\begin{namedtheorem}[6.5.7]{Remarques}
\textnormal{(i)} Considérons en particulier le cas où $S=X=Y$,
$\mathscr P_\bullet$ et $\mathscr Q_\bullet$ étant donc deux complexes de
$\mathscr O_S$-Modules quasi-cohérents; alors les
$\mathscr{T}\!or_n^S(\mathscr P_\bullet,\mathscr Q_\bullet)$ sont des
$\mathscr O_S$-Modules quasi-cohérents; en outre, pour tout point $z\in S$, il
résulte de (6.5.6) que l'on a un isomorphisme canonique
\[
  \bigl(\mathscr{T}\!or_n^S
    (\mathscr P_\bullet,\mathscr Q_\bullet)\bigr)_z
  \xrightarrow{\sim}
  \mathbf{Tor}_n^{\mathscr O_z}
  \bigl((\mathscr P_\bullet)_z,(\mathscr Q_\bullet)_z\bigr)
  \tag{6.5.7.1}\label{III.6.5.7.1.fr}
\]
car la question est locale et on est ramené au cas des modules, en vertu de
(6.4.1.1).

\textnormal{(ii)} On peut généraliser la définition des hypertor au cas de
deux complexes de $\mathscr O_X$-Modules $\mathscr P_\bullet$,
$\mathscr Q_\bullet$ sur un même espace annelé $(X,\mathscr O_X)$; pour tout
ouvert $U$ de $X$, posons en effet
$A(U)=\Gamma(U,\mathscr O_X)$,
$P_\bullet(U)=\Gamma(U,\mathscr P_\bullet)$,
$Q_\bullet(U)=\Gamma(U,\mathscr Q_\bullet)$; les $A(U)$-modules
\[
  \mathbf{Tor}_n^{A(U)}(P_\bullet(U),Q_\bullet(U))
\]
forment alors un préfaisceau sur $X$, et l'on désigne par
$\mathscr{T}\!or_n^X(\mathscr P_\bullet,\mathscr Q_\bullet)$ le
$\mathscr O_X$-Module associé à ce préfaisceau. Lorsque $X$ est un préschéma,
il résulte de (6.3.12) que ce $\mathscr O_X$-Module est canoniquement isomorphe
à l'hypertor défini ci-dessus. Nous ne développerons pas davantage cette
généralisation.
\end{namedtheorem}

\begin{proposition}[6.5.8]
Soient $X$, $Y$ deux $S$-préschémas, $\mathscr F$ (resp. $\mathscr G$) un
$\mathscr O_X$-Module (resp. un $\mathscr O_Y$-Module) quasi-cohérent. Si
$\mathscr F$ ou $\mathscr G$ est $S$-plat, on a
$\mathscr{T}\!or_n^S(\mathscr F,\mathscr G)=0$ pour $n\neq0$.
\end{proposition}

\begin{proof}
La question étant locale sur $X$ et $Y$, on peut supposer $X$, $Y$ et $S$
affines, d'anneaux respectifs $B$, $C$, $A$, et
$\mathscr F=\widetilde M$, $\mathscr G=\widetilde N$, $M$ (resp. $N$) étant
un $B$-module (resp. un $C$-module). Supposons par exemple que $\mathscr F$
soit $S$-plat, ce qui signifie que pour tout $s\in S$, $M_s$ est un
$A_s$-module plat ($0_{\mathrm I}$, 6.7.1); par suite $M$ est un $A$-module plat
($0_{\mathrm I}$, 6.3.3), et l'on sait que
$\operatorname{Tor}_n^A(M,N)=0$ pour $n>0$ et pour tout $C$-module $N$
($0_{\mathrm I}$, 6.1.1), d'où la conclusion par (6.4.1.1).
\end{proof}

\begin{namedtheorem}[6.5.9]{Corollaire}
Soient $X$, $Y$ deux $S$-préschémas, $\mathscr P_\bullet$ (resp.
$\mathscr Q_\bullet$) un complexe de

